% The textbook design of the full book, shared by every edition: % The textbook design of the full book, shared by every edition: % The textbook design of the full book, shared by every edition: % The textbook design of the full book, shared by every edition: \input{figures/booklook} at the end of the preamble.
% Type: Libertinus Serif for text, Libertinus Math for formulas, Source Sans 3 for headings, captions, figure labels
% and running heads. Chapter openers: a full-width band with the number at the outer edge, the title in sans below,
% and the "In this chapter" box (environment chaptergoals, defined plainly in each main.tex for the web build).
% Propositions: tinted panels with an accent rule; exercises: "Exercise 4.3 ** Title" in sans.
% An edition may define before the \input:
%   \chapterlabel         the word in the band (default: \@chapapp, i.e. Chapter / Appendix from babel;
%                         \booklookapponly = the word only on appendices, for languages that put the number inside it);
%   \booklooknohead       keep the edition's own running heads (Arabic: right-to-left page numbers);
%   \booklookbarside      the side of the accent rule on key statements (default west; east for right-to-left);
%   \booklookbabelfonts   set the Latin faces with \babelfont[english] (Arabic) instead of \setmainfont.
\makeatletter
% ---------- type
\RequirePackage{unicode-math}
\ifdefined\booklookbabelfonts
  \babelfont[english]{rm}{Libertinus Serif}
  \babelfont[english]{sf}[Scale=MatchLowercase,BoldFont=* Semibold]{Source Sans 3}
\else
  \setmainfont{Libertinus Serif}
  \setsansfont{Source Sans 3}[Scale=MatchLowercase,BoldFont=* Semibold]
\fi
\setmathfont{Libertinus Math}
% a loose line rather than one running into the margin (short lines with long code names, e.g. exercise titles)
\setlength{\emergencystretch}{6em}

% ---------- colour
\definecolor{accent}{rgb}{0.11,0.30,0.53}
\colorlet{accentdark}{accent!75!black}
\colorlet{tint}{accent!6}
\colorlet{inkgray}{black!62}
\hypersetup{linkcolor=accent,citecolor=accent,urlcolor=accent}
\providecommand{\chapterlabel}{\@chapapp}
\providecommand{\booklookbarside}{west}
\def\booklook@app{\appendixname}
\newcommand{\booklookapponly}{\ifx\@chapapp\booklook@app\appendixname\fi}

% ---------- chapter opener: a full-width band, the number at the outer edge, the title below in sans
% the band: an overlay anchored to the page (needs the positions of the previous run; the Makefile runs three passes)
\def\booklook@band#1{%
  \noindent\begin{tikzpicture}[remember picture,overlay]
    \fill[accent] (current page.north west) rectangle ([yshift=-1.45in]current page.north east);
    \fill[accentdark] ([yshift=-1.45in]current page.north west) rectangle ([yshift=-1.51in]current page.north east);
    #1
  \end{tikzpicture}\par\nointerlineskip}
\def\booklook@title#1{%
  \vspace*{\dimexpr0.6in-\topskip\relax}%
  {\parindent\z@\raggedright\sffamily\bfseries\fontsize{25}{30}\selectfont\color{accentdark}#1\par\nobreak}%
  \vskip 20\p@}
\def\@makechapterhead#1{%
  \ifnum\c@secnumdepth>\m@ne\if@mainmatter
    \booklook@band{%
      \ifodd\c@page
        \node[anchor=base east,text=white,inner sep=0pt,font=\sffamily\bfseries\fontsize{80}{80}\selectfont]
          at ([xshift=-0.55in,yshift=-1.22in]current page.north east) {\thechapter};
        \node[anchor=base west,text=white!80!accent,inner sep=0pt,font=\sffamily\large]
          at ([xshift=\Gm@lmargin,yshift=-1.22in]current page.north west) {\MakeUppercase{\chapterlabel}};
      \else
        \node[anchor=base west,text=white,inner sep=0pt,font=\sffamily\bfseries\fontsize{80}{80}\selectfont]
          at ([xshift=0.55in,yshift=-1.22in]current page.north west) {\thechapter};
        \node[anchor=base east,text=white!80!accent,inner sep=0pt,font=\sffamily\large]
          at ([xshift=-\Gm@rmargin,yshift=-1.22in]current page.north east) {\MakeUppercase{\chapterlabel}};
      \fi}%
  \else\booklook@band{}\fi\else\booklook@band{}\fi
  \booklook@title{#1}}
\def\@makeschapterhead#1{\booklook@band{}\booklook@title{#1}}
% "In this chapter": the chapter's goals under its title
\renewenvironment{chaptergoals}{%
  \begin{tcolorbox}[enhanced,colback=tint,colframe=tint,arc=0pt,boxrule=0pt,left=10pt,right=10pt,top=6pt,bottom=6pt,
    fontupper=\sffamily\small,before skip=0pt,after skip=14pt]
  {\color{accent}\bfseries\MakeUppercase{\chaptergoalsname}}\par\smallskip
  \begin{itemize}\setlength{\itemsep}{1pt}}{\end{itemize}\end{tcolorbox}}

% ---------- headings in sans, numbers in the accent colour
\def\@seccntformat#1{{\color{accent}\csname the#1\endcsname}\hskip 0.7em}
\renewcommand\section{\@startsection{section}{1}{\z@}{-3.5ex \@plus -1ex \@minus -.2ex}{2ex \@plus.2ex}{\normalfont\sffamily\bfseries\Large\color{accentdark}\raggedright}}
\renewcommand\subsection{\@startsection{subsection}{2}{\z@}{-3.25ex\@plus -1ex \@minus -.2ex}{1.5ex \@plus .2ex}{\normalfont\sffamily\bfseries\large\color{accentdark}\raggedright}}
\renewcommand\subsubsection{\@startsection{subsubsection}{3}{\z@}{-3.25ex\@plus -1ex \@minus -.2ex}{1.5ex \@plus .2ex}{\normalfont\sffamily\bfseries\normalsize\color{accentdark}}}

% ---------- running heads in sans: the chapter and section marks, a thin accent rule, the page number outside
\ifdefined\booklooknohead\else
\def\booklook@rule{\makebox[\z@][l]{\color{accent}\rule[-5\p@]{\textwidth}{0.4\p@}}}
\def\ps@headings{%
  \let\@oddfoot\@empty\let\@evenfoot\@empty
  \def\@evenhead{\sffamily\booklook@rule{\bfseries\color{accent}\thepage}\hfil{\color{inkgray}\leftmark}}%
  \def\@oddhead{\sffamily\booklook@rule{\color{inkgray}\rightmark}\hfil{\bfseries\color{accent}\thepage}}%
  \let\@mkboth\markboth}
\def\ps@plain{\let\@oddhead\@empty\let\@evenhead\@empty
  \def\@oddfoot{\hfil\sffamily\color{inkgray}\thepage\hfil}\let\@evenfoot\@oddfoot}
\pagestyle{headings}
\fi

% ---------- figures and tables in sans (labels inside the drawings, table text, captions); captions small
\g@addto@macro\@floatboxreset{\sffamily}
\def\fnum@figure{\textcolor{accent}{\bfseries\figurename\nobreakspace\thefigure}}
\def\fnum@table{\textcolor{accent}{\bfseries\tablename\nobreakspace\thetable}}
\AtBeginDocument{\let\booklook@makecaption\@makecaption
  \long\def\@makecaption#1#2{\small\booklook@makecaption{#1}{#2}}}

% ---------- key statements: a tinted panel with an accent rule on the leading edge
\def\booklook@west{west}
\ifx\booklookbarside\booklook@west\def\booklook@rulekey{leftrule}\else\def\booklook@rulekey{rightrule}\fi
\tcbset{booklook rule/.style/.expanded={\booklook@rulekey=2.5pt}}
\renewtcolorbox{keyblock}{enhanced,breakable,colback=tint,colframe=accent,boxrule=0pt,arc=0pt,booklook rule,
  left=9pt,right=9pt,top=5pt,bottom=5pt,boxsep=0pt}
% theorem heads (the names come from each edition's \newtheorem): "Proposition 4.1 · Title." in sans accent,
% exercises "Exercise 4.3 ** Title" with the difficulty stars; amsthm looks the style up when a theorem is used
\newtheoremstyle{plain}{\topsep}{\topsep}{\itshape}{}{\sffamily\bfseries\color{accent}}{.}{0.5em}
  {\thmname{#1}\thmnumber{~#2}\thmnote{\ \textperiodcentered\ #3}}
\newtheoremstyle{definition}{\topsep}{\topsep}{\normalfont}{}{\sffamily\bfseries\color{accentdark}}{}{0.6em}
  {\thmname{#1}\thmnumber{~#2}\thmnote{\ #3}}
\makeatother
 at the end of the preamble.
% Type: Libertinus Serif for text, Libertinus Math for formulas, Source Sans 3 for headings, captions, figure labels
% and running heads. Chapter openers: a full-width band with the number at the outer edge, the title in sans below,
% and the "In this chapter" box (environment chaptergoals, defined plainly in each main.tex for the web build).
% Propositions: tinted panels with an accent rule; exercises: "Exercise 4.3 ** Title" in sans.
% An edition may define before the \input:
%   \chapterlabel         the word in the band (default: \@chapapp, i.e. Chapter / Appendix from babel;
%                         \booklookapponly = the word only on appendices, for languages that put the number inside it);
%   \booklooknohead       keep the edition's own running heads (Arabic: right-to-left page numbers);
%   \booklookbarside      the side of the accent rule on key statements (default west; east for right-to-left);
%   \booklookbabelfonts   set the Latin faces with \babelfont[english] (Arabic) instead of \setmainfont.
\makeatletter
% ---------- type
\RequirePackage{unicode-math}
\ifdefined\booklookbabelfonts
  \babelfont[english]{rm}{Libertinus Serif}
  \babelfont[english]{sf}[Scale=MatchLowercase,BoldFont=* Semibold]{Source Sans 3}
\else
  \setmainfont{Libertinus Serif}
  \setsansfont{Source Sans 3}[Scale=MatchLowercase,BoldFont=* Semibold]
\fi
\setmathfont{Libertinus Math}
% a loose line rather than one running into the margin (short lines with long code names, e.g. exercise titles)
\setlength{\emergencystretch}{6em}

% ---------- colour
\definecolor{accent}{rgb}{0.11,0.30,0.53}
\colorlet{accentdark}{accent!75!black}
\colorlet{tint}{accent!6}
\colorlet{inkgray}{black!62}
\hypersetup{linkcolor=accent,citecolor=accent,urlcolor=accent}
\providecommand{\chapterlabel}{\@chapapp}
\providecommand{\booklookbarside}{west}
\def\booklook@app{\appendixname}
\newcommand{\booklookapponly}{\ifx\@chapapp\booklook@app\appendixname\fi}

% ---------- chapter opener: a full-width band, the number at the outer edge, the title below in sans
% the band: an overlay anchored to the page (needs the positions of the previous run; the Makefile runs three passes)
\def\booklook@band#1{%
  \noindent\begin{tikzpicture}[remember picture,overlay]
    \fill[accent] (current page.north west) rectangle ([yshift=-1.45in]current page.north east);
    \fill[accentdark] ([yshift=-1.45in]current page.north west) rectangle ([yshift=-1.51in]current page.north east);
    #1
  \end{tikzpicture}\par\nointerlineskip}
\def\booklook@title#1{%
  \vspace*{\dimexpr0.6in-\topskip\relax}%
  {\parindent\z@\raggedright\sffamily\bfseries\fontsize{25}{30}\selectfont\color{accentdark}#1\par\nobreak}%
  \vskip 20\p@}
\def\@makechapterhead#1{%
  \ifnum\c@secnumdepth>\m@ne\if@mainmatter
    \booklook@band{%
      \ifodd\c@page
        \node[anchor=base east,text=white,inner sep=0pt,font=\sffamily\bfseries\fontsize{80}{80}\selectfont]
          at ([xshift=-0.55in,yshift=-1.22in]current page.north east) {\thechapter};
        \node[anchor=base west,text=white!80!accent,inner sep=0pt,font=\sffamily\large]
          at ([xshift=\Gm@lmargin,yshift=-1.22in]current page.north west) {\MakeUppercase{\chapterlabel}};
      \else
        \node[anchor=base west,text=white,inner sep=0pt,font=\sffamily\bfseries\fontsize{80}{80}\selectfont]
          at ([xshift=0.55in,yshift=-1.22in]current page.north west) {\thechapter};
        \node[anchor=base east,text=white!80!accent,inner sep=0pt,font=\sffamily\large]
          at ([xshift=-\Gm@rmargin,yshift=-1.22in]current page.north east) {\MakeUppercase{\chapterlabel}};
      \fi}%
  \else\booklook@band{}\fi\else\booklook@band{}\fi
  \booklook@title{#1}}
\def\@makeschapterhead#1{\booklook@band{}\booklook@title{#1}}
% "In this chapter": the chapter's goals under its title
\renewenvironment{chaptergoals}{%
  \begin{tcolorbox}[enhanced,colback=tint,colframe=tint,arc=0pt,boxrule=0pt,left=10pt,right=10pt,top=6pt,bottom=6pt,
    fontupper=\sffamily\small,before skip=0pt,after skip=14pt]
  {\color{accent}\bfseries\MakeUppercase{\chaptergoalsname}}\par\smallskip
  \begin{itemize}\setlength{\itemsep}{1pt}}{\end{itemize}\end{tcolorbox}}

% ---------- headings in sans, numbers in the accent colour
\def\@seccntformat#1{{\color{accent}\csname the#1\endcsname}\hskip 0.7em}
\renewcommand\section{\@startsection{section}{1}{\z@}{-3.5ex \@plus -1ex \@minus -.2ex}{2ex \@plus.2ex}{\normalfont\sffamily\bfseries\Large\color{accentdark}\raggedright}}
\renewcommand\subsection{\@startsection{subsection}{2}{\z@}{-3.25ex\@plus -1ex \@minus -.2ex}{1.5ex \@plus .2ex}{\normalfont\sffamily\bfseries\large\color{accentdark}\raggedright}}
\renewcommand\subsubsection{\@startsection{subsubsection}{3}{\z@}{-3.25ex\@plus -1ex \@minus -.2ex}{1.5ex \@plus .2ex}{\normalfont\sffamily\bfseries\normalsize\color{accentdark}}}

% ---------- running heads in sans: the chapter and section marks, a thin accent rule, the page number outside
\ifdefined\booklooknohead\else
\def\booklook@rule{\makebox[\z@][l]{\color{accent}\rule[-5\p@]{\textwidth}{0.4\p@}}}
\def\ps@headings{%
  \let\@oddfoot\@empty\let\@evenfoot\@empty
  \def\@evenhead{\sffamily\booklook@rule{\bfseries\color{accent}\thepage}\hfil{\color{inkgray}\leftmark}}%
  \def\@oddhead{\sffamily\booklook@rule{\color{inkgray}\rightmark}\hfil{\bfseries\color{accent}\thepage}}%
  \let\@mkboth\markboth}
\def\ps@plain{\let\@oddhead\@empty\let\@evenhead\@empty
  \def\@oddfoot{\hfil\sffamily\color{inkgray}\thepage\hfil}\let\@evenfoot\@oddfoot}
\pagestyle{headings}
\fi

% ---------- figures and tables in sans (labels inside the drawings, table text, captions); captions small
\g@addto@macro\@floatboxreset{\sffamily}
\def\fnum@figure{\textcolor{accent}{\bfseries\figurename\nobreakspace\thefigure}}
\def\fnum@table{\textcolor{accent}{\bfseries\tablename\nobreakspace\thetable}}
\AtBeginDocument{\let\booklook@makecaption\@makecaption
  \long\def\@makecaption#1#2{\small\booklook@makecaption{#1}{#2}}}

% ---------- key statements: a tinted panel with an accent rule on the leading edge
\def\booklook@west{west}
\ifx\booklookbarside\booklook@west\def\booklook@rulekey{leftrule}\else\def\booklook@rulekey{rightrule}\fi
\tcbset{booklook rule/.style/.expanded={\booklook@rulekey=2.5pt}}
\renewtcolorbox{keyblock}{enhanced,breakable,colback=tint,colframe=accent,boxrule=0pt,arc=0pt,booklook rule,
  left=9pt,right=9pt,top=5pt,bottom=5pt,boxsep=0pt}
% theorem heads (the names come from each edition's \newtheorem): "Proposition 4.1 · Title." in sans accent,
% exercises "Exercise 4.3 ** Title" with the difficulty stars; amsthm looks the style up when a theorem is used
\newtheoremstyle{plain}{\topsep}{\topsep}{\itshape}{}{\sffamily\bfseries\color{accent}}{.}{0.5em}
  {\thmname{#1}\thmnumber{~#2}\thmnote{\ \textperiodcentered\ #3}}
\newtheoremstyle{definition}{\topsep}{\topsep}{\normalfont}{}{\sffamily\bfseries\color{accentdark}}{}{0.6em}
  {\thmname{#1}\thmnumber{~#2}\thmnote{\ #3}}
\makeatother
 at the end of the preamble.
% Type: Libertinus Serif for text, Libertinus Math for formulas, Source Sans 3 for headings, captions, figure labels
% and running heads. Chapter openers: a full-width band with the number at the outer edge, the title in sans below,
% and the "In this chapter" box (environment chaptergoals, defined plainly in each main.tex for the web build).
% Propositions: tinted panels with an accent rule; exercises: "Exercise 4.3 ** Title" in sans.
% An edition may define before the \input:
%   \chapterlabel         the word in the band (default: \@chapapp, i.e. Chapter / Appendix from babel;
%                         \booklookapponly = the word only on appendices, for languages that put the number inside it);
%   \booklooknohead       keep the edition's own running heads (Arabic: right-to-left page numbers);
%   \booklookbarside      the side of the accent rule on key statements (default west; east for right-to-left);
%   \booklookbabelfonts   set the Latin faces with \babelfont[english] (Arabic) instead of \setmainfont.
\makeatletter
% ---------- type
\RequirePackage{unicode-math}
\ifdefined\booklookbabelfonts
  \babelfont[english]{rm}{Libertinus Serif}
  \babelfont[english]{sf}[Scale=MatchLowercase,BoldFont=* Semibold]{Source Sans 3}
\else
  \setmainfont{Libertinus Serif}
  \setsansfont{Source Sans 3}[Scale=MatchLowercase,BoldFont=* Semibold]
\fi
\setmathfont{Libertinus Math}
% a loose line rather than one running into the margin (short lines with long code names, e.g. exercise titles)
\setlength{\emergencystretch}{6em}

% ---------- colour
\definecolor{accent}{rgb}{0.11,0.30,0.53}
\colorlet{accentdark}{accent!75!black}
\colorlet{tint}{accent!6}
\colorlet{inkgray}{black!62}
\hypersetup{linkcolor=accent,citecolor=accent,urlcolor=accent}
\providecommand{\chapterlabel}{\@chapapp}
\providecommand{\booklookbarside}{west}
\def\booklook@app{\appendixname}
\newcommand{\booklookapponly}{\ifx\@chapapp\booklook@app\appendixname\fi}

% ---------- chapter opener: a full-width band, the number at the outer edge, the title below in sans
% the band: an overlay anchored to the page (needs the positions of the previous run; the Makefile runs three passes)
\def\booklook@band#1{%
  \noindent\begin{tikzpicture}[remember picture,overlay]
    \fill[accent] (current page.north west) rectangle ([yshift=-1.45in]current page.north east);
    \fill[accentdark] ([yshift=-1.45in]current page.north west) rectangle ([yshift=-1.51in]current page.north east);
    #1
  \end{tikzpicture}\par\nointerlineskip}
\def\booklook@title#1{%
  \vspace*{\dimexpr0.6in-\topskip\relax}%
  {\parindent\z@\raggedright\sffamily\bfseries\fontsize{25}{30}\selectfont\color{accentdark}#1\par\nobreak}%
  \vskip 20\p@}
\def\@makechapterhead#1{%
  \ifnum\c@secnumdepth>\m@ne\if@mainmatter
    \booklook@band{%
      \ifodd\c@page
        \node[anchor=base east,text=white,inner sep=0pt,font=\sffamily\bfseries\fontsize{80}{80}\selectfont]
          at ([xshift=-0.55in,yshift=-1.22in]current page.north east) {\thechapter};
        \node[anchor=base west,text=white!80!accent,inner sep=0pt,font=\sffamily\large]
          at ([xshift=\Gm@lmargin,yshift=-1.22in]current page.north west) {\MakeUppercase{\chapterlabel}};
      \else
        \node[anchor=base west,text=white,inner sep=0pt,font=\sffamily\bfseries\fontsize{80}{80}\selectfont]
          at ([xshift=0.55in,yshift=-1.22in]current page.north west) {\thechapter};
        \node[anchor=base east,text=white!80!accent,inner sep=0pt,font=\sffamily\large]
          at ([xshift=-\Gm@rmargin,yshift=-1.22in]current page.north east) {\MakeUppercase{\chapterlabel}};
      \fi}%
  \else\booklook@band{}\fi\else\booklook@band{}\fi
  \booklook@title{#1}}
\def\@makeschapterhead#1{\booklook@band{}\booklook@title{#1}}
% "In this chapter": the chapter's goals under its title
\renewenvironment{chaptergoals}{%
  \begin{tcolorbox}[enhanced,colback=tint,colframe=tint,arc=0pt,boxrule=0pt,left=10pt,right=10pt,top=6pt,bottom=6pt,
    fontupper=\sffamily\small,before skip=0pt,after skip=14pt]
  {\color{accent}\bfseries\MakeUppercase{\chaptergoalsname}}\par\smallskip
  \begin{itemize}\setlength{\itemsep}{1pt}}{\end{itemize}\end{tcolorbox}}

% ---------- headings in sans, numbers in the accent colour
\def\@seccntformat#1{{\color{accent}\csname the#1\endcsname}\hskip 0.7em}
\renewcommand\section{\@startsection{section}{1}{\z@}{-3.5ex \@plus -1ex \@minus -.2ex}{2ex \@plus.2ex}{\normalfont\sffamily\bfseries\Large\color{accentdark}\raggedright}}
\renewcommand\subsection{\@startsection{subsection}{2}{\z@}{-3.25ex\@plus -1ex \@minus -.2ex}{1.5ex \@plus .2ex}{\normalfont\sffamily\bfseries\large\color{accentdark}\raggedright}}
\renewcommand\subsubsection{\@startsection{subsubsection}{3}{\z@}{-3.25ex\@plus -1ex \@minus -.2ex}{1.5ex \@plus .2ex}{\normalfont\sffamily\bfseries\normalsize\color{accentdark}}}

% ---------- running heads in sans: the chapter and section marks, a thin accent rule, the page number outside
\ifdefined\booklooknohead\else
\def\booklook@rule{\makebox[\z@][l]{\color{accent}\rule[-5\p@]{\textwidth}{0.4\p@}}}
\def\ps@headings{%
  \let\@oddfoot\@empty\let\@evenfoot\@empty
  \def\@evenhead{\sffamily\booklook@rule{\bfseries\color{accent}\thepage}\hfil{\color{inkgray}\leftmark}}%
  \def\@oddhead{\sffamily\booklook@rule{\color{inkgray}\rightmark}\hfil{\bfseries\color{accent}\thepage}}%
  \let\@mkboth\markboth}
\def\ps@plain{\let\@oddhead\@empty\let\@evenhead\@empty
  \def\@oddfoot{\hfil\sffamily\color{inkgray}\thepage\hfil}\let\@evenfoot\@oddfoot}
\pagestyle{headings}
\fi

% ---------- figures and tables in sans (labels inside the drawings, table text, captions); captions small
\g@addto@macro\@floatboxreset{\sffamily}
\def\fnum@figure{\textcolor{accent}{\bfseries\figurename\nobreakspace\thefigure}}
\def\fnum@table{\textcolor{accent}{\bfseries\tablename\nobreakspace\thetable}}
\AtBeginDocument{\let\booklook@makecaption\@makecaption
  \long\def\@makecaption#1#2{\small\booklook@makecaption{#1}{#2}}}

% ---------- key statements: a tinted panel with an accent rule on the leading edge
\def\booklook@west{west}
\ifx\booklookbarside\booklook@west\def\booklook@rulekey{leftrule}\else\def\booklook@rulekey{rightrule}\fi
\tcbset{booklook rule/.style/.expanded={\booklook@rulekey=2.5pt}}
\renewtcolorbox{keyblock}{enhanced,breakable,colback=tint,colframe=accent,boxrule=0pt,arc=0pt,booklook rule,
  left=9pt,right=9pt,top=5pt,bottom=5pt,boxsep=0pt}
% theorem heads (the names come from each edition's \newtheorem): "Proposition 4.1 · Title." in sans accent,
% exercises "Exercise 4.3 ** Title" with the difficulty stars; amsthm looks the style up when a theorem is used
\newtheoremstyle{plain}{\topsep}{\topsep}{\itshape}{}{\sffamily\bfseries\color{accent}}{.}{0.5em}
  {\thmname{#1}\thmnumber{~#2}\thmnote{\ \textperiodcentered\ #3}}
\newtheoremstyle{definition}{\topsep}{\topsep}{\normalfont}{}{\sffamily\bfseries\color{accentdark}}{}{0.6em}
  {\thmname{#1}\thmnumber{~#2}\thmnote{\ #3}}
\makeatother
 at the end of the preamble.
% Type: Libertinus Serif for text, Libertinus Math for formulas, Source Sans 3 for headings, captions, figure labels
% and running heads. Chapter openers: a full-width band with the number at the outer edge, the title in sans below,
% and the "In this chapter" box (environment chaptergoals, defined plainly in each main.tex for the web build).
% Propositions: tinted panels with an accent rule; exercises: "Exercise 4.3 ** Title" in sans.
% An edition may define before the \input:
%   \chapterlabel         the word in the band (default: \@chapapp, i.e. Chapter / Appendix from babel;
%                         \booklookapponly = the word only on appendices, for languages that put the number inside it);
%   \booklooknohead       keep the edition's own running heads (Arabic: right-to-left page numbers);
%   \booklookbarside      the side of the accent rule on key statements (default west; east for right-to-left);
%   \booklookbabelfonts   set the Latin faces with \babelfont[english] (Arabic) instead of \setmainfont.
\makeatletter
% ---------- type
\RequirePackage{unicode-math}
\ifdefined\booklookbabelfonts
  \babelfont[english]{rm}{Libertinus Serif}
  \babelfont[english]{sf}[Scale=MatchLowercase,BoldFont=* Semibold]{Source Sans 3}
\else
  \setmainfont{Libertinus Serif}
  \setsansfont{Source Sans 3}[Scale=MatchLowercase,BoldFont=* Semibold]
\fi
\setmathfont{Libertinus Math}
% a loose line rather than one running into the margin (short lines with long code names, e.g. exercise titles)
\setlength{\emergencystretch}{6em}

% ---------- colour
\definecolor{accent}{rgb}{0.11,0.30,0.53}
\colorlet{accentdark}{accent!75!black}
\colorlet{tint}{accent!6}
\colorlet{inkgray}{black!62}
\hypersetup{linkcolor=accent,citecolor=accent,urlcolor=accent}
\providecommand{\chapterlabel}{\@chapapp}
\providecommand{\booklookbarside}{west}
\def\booklook@app{\appendixname}
\newcommand{\booklookapponly}{\ifx\@chapapp\booklook@app\appendixname\fi}

% ---------- chapter opener: a full-width band, the number at the outer edge, the title below in sans
% the band: an overlay anchored to the page (needs the positions of the previous run; the Makefile runs three passes)
\def\booklook@band#1{%
  \noindent\begin{tikzpicture}[remember picture,overlay]
    \fill[accent] (current page.north west) rectangle ([yshift=-1.45in]current page.north east);
    \fill[accentdark] ([yshift=-1.45in]current page.north west) rectangle ([yshift=-1.51in]current page.north east);
    #1
  \end{tikzpicture}\par\nointerlineskip}
\def\booklook@title#1{%
  \vspace*{\dimexpr0.6in-\topskip\relax}%
  {\parindent\z@\raggedright\sffamily\bfseries\fontsize{25}{30}\selectfont\color{accentdark}#1\par\nobreak}%
  \vskip 20\p@}
\def\@makechapterhead#1{%
  \ifnum\c@secnumdepth>\m@ne\if@mainmatter
    \booklook@band{%
      \ifodd\c@page
        \node[anchor=base east,text=white,inner sep=0pt,font=\sffamily\bfseries\fontsize{80}{80}\selectfont]
          at ([xshift=-0.55in,yshift=-1.22in]current page.north east) {\thechapter};
        \node[anchor=base west,text=white!80!accent,inner sep=0pt,font=\sffamily\large]
          at ([xshift=\Gm@lmargin,yshift=-1.22in]current page.north west) {\MakeUppercase{\chapterlabel}};
      \else
        \node[anchor=base west,text=white,inner sep=0pt,font=\sffamily\bfseries\fontsize{80}{80}\selectfont]
          at ([xshift=0.55in,yshift=-1.22in]current page.north west) {\thechapter};
        \node[anchor=base east,text=white!80!accent,inner sep=0pt,font=\sffamily\large]
          at ([xshift=-\Gm@rmargin,yshift=-1.22in]current page.north east) {\MakeUppercase{\chapterlabel}};
      \fi}%
  \else\booklook@band{}\fi\else\booklook@band{}\fi
  \booklook@title{#1}}
\def\@makeschapterhead#1{\booklook@band{}\booklook@title{#1}}
% "In this chapter": the chapter's goals under its title
\renewenvironment{chaptergoals}{%
  \begin{tcolorbox}[enhanced,colback=tint,colframe=tint,arc=0pt,boxrule=0pt,left=10pt,right=10pt,top=6pt,bottom=6pt,
    fontupper=\sffamily\small,before skip=0pt,after skip=14pt]
  {\color{accent}\bfseries\MakeUppercase{\chaptergoalsname}}\par\smallskip
  \begin{itemize}\setlength{\itemsep}{1pt}}{\end{itemize}\end{tcolorbox}}

% ---------- headings in sans, numbers in the accent colour
\def\@seccntformat#1{{\color{accent}\csname the#1\endcsname}\hskip 0.7em}
\renewcommand\section{\@startsection{section}{1}{\z@}{-3.5ex \@plus -1ex \@minus -.2ex}{2ex \@plus.2ex}{\normalfont\sffamily\bfseries\Large\color{accentdark}\raggedright}}
\renewcommand\subsection{\@startsection{subsection}{2}{\z@}{-3.25ex\@plus -1ex \@minus -.2ex}{1.5ex \@plus .2ex}{\normalfont\sffamily\bfseries\large\color{accentdark}\raggedright}}
\renewcommand\subsubsection{\@startsection{subsubsection}{3}{\z@}{-3.25ex\@plus -1ex \@minus -.2ex}{1.5ex \@plus .2ex}{\normalfont\sffamily\bfseries\normalsize\color{accentdark}}}

% ---------- running heads in sans: the chapter and section marks, a thin accent rule, the page number outside
\ifdefined\booklooknohead\else
\def\booklook@rule{\makebox[\z@][l]{\color{accent}\rule[-5\p@]{\textwidth}{0.4\p@}}}
\def\ps@headings{%
  \let\@oddfoot\@empty\let\@evenfoot\@empty
  \def\@evenhead{\sffamily\booklook@rule{\bfseries\color{accent}\thepage}\hfil{\color{inkgray}\leftmark}}%
  \def\@oddhead{\sffamily\booklook@rule{\color{inkgray}\rightmark}\hfil{\bfseries\color{accent}\thepage}}%
  \let\@mkboth\markboth}
\def\ps@plain{\let\@oddhead\@empty\let\@evenhead\@empty
  \def\@oddfoot{\hfil\sffamily\color{inkgray}\thepage\hfil}\let\@evenfoot\@oddfoot}
\pagestyle{headings}
\fi

% ---------- figures and tables in sans (labels inside the drawings, table text, captions); captions small
\g@addto@macro\@floatboxreset{\sffamily}
\def\fnum@figure{\textcolor{accent}{\bfseries\figurename\nobreakspace\thefigure}}
\def\fnum@table{\textcolor{accent}{\bfseries\tablename\nobreakspace\thetable}}
\AtBeginDocument{\let\booklook@makecaption\@makecaption
  \long\def\@makecaption#1#2{\small\booklook@makecaption{#1}{#2}}}

% ---------- key statements: a tinted panel with an accent rule on the leading edge
\def\booklook@west{west}
\ifx\booklookbarside\booklook@west\def\booklook@rulekey{leftrule}\else\def\booklook@rulekey{rightrule}\fi
\tcbset{booklook rule/.style/.expanded={\booklook@rulekey=2.5pt}}
\renewtcolorbox{keyblock}{enhanced,breakable,colback=tint,colframe=accent,boxrule=0pt,arc=0pt,booklook rule,
  left=9pt,right=9pt,top=5pt,bottom=5pt,boxsep=0pt}
% theorem heads (the names come from each edition's \newtheorem): "Proposition 4.1 · Title." in sans accent,
% exercises "Exercise 4.3 ** Title" with the difficulty stars; amsthm looks the style up when a theorem is used
\newtheoremstyle{plain}{\topsep}{\topsep}{\itshape}{}{\sffamily\bfseries\color{accent}}{.}{0.5em}
  {\thmname{#1}\thmnumber{~#2}\thmnote{\ \textperiodcentered\ #3}}
\newtheoremstyle{definition}{\topsep}{\topsep}{\normalfont}{}{\sffamily\bfseries\color{accentdark}}{}{0.6em}
  {\thmname{#1}\thmnumber{~#2}\thmnote{\ #3}}
\makeatother
